\begin{proposition}[DP39: what harmonicity removes, and what it does not]\label{DP39}
The following two assertions distinguish full-series denominators from selected theta-type jets.
\begin{enumerate}[leftmargin=*,label=(\alph*)]
\item There are integral formal-biextension data, two canonical gauges, and a primitive germ $f=1$ for which one normalized series has unbounded denominators, while all the specified harmonic values agree. In fact the agreement holds for every germ $f$ whenever the polynomial $P$ is killed by the specified second-order contraction.
\item For odd $p$ and every $h\ge2$, two copies of the actual primitive pair have an integral polynomial satisfying the summed harmonic, row-invariance and fixed-row-degree conditions defining \path{HaO}, whose normalized kernel value is a unit. The same invariant polynomial takes values in $pR$ on every integral function on that product. Therefore those algebraic harmonic conditions alone do not transfer the bare integral-function divisibility through the primitive exponential.
\end{enumerate}
Neither assertion constructs a new global mixed-type theta correspondence or asserts that the particular germ $1$ is a global theta section.
\end{proposition}
\begin{proof}
For a polynomial $P$ and a formal series $f$ over a characteristic-zero field, write $\langle P,f\rangle=P(\partial)f(0)$. For every polynomial $Q$,
\begin{equation}\label{eq:fischer-ideal}
 \langle P,Qf\rangle=\langle Q(\partial)P,f\rangle.
\end{equation}
For monomials $P=Z^\alpha,Q=Z^\beta,f=Z^\gamma$, both sides vanish unless $\alpha=\beta+\gamma$, when both are $\alpha!$. Linearity and the finite jet read by $P$ prove the identity for formal series. If $Q_j(\partial)P=0$ for every $j$, the pairing kills the formal ideal $(Q_j)$. If $\eta$ has zero constant term and belongs to that ideal, then $\exp\eta-1=\eta\sum_{n\ge1}\eta^{n-1}/n!$ also belongs to it. Hence
\[
 \langle P,f\exp\eta\rangle=\langle P,f\rangle.
\]

For (a), let $F=\widehat\Ga^2$ over $\Z_p$, let $O=\Z\times\Z$ act diagonally, and let its involution exchange the two factors. Take trivial addition, symmetry and adjoint cocycles. Both $h_0=0$ and
\[
 h_1(X,Y)=X_1Y_2+X_2Y_1
\]
are canonical gauges: their axes vanish, their addition coboundaries vanish by biadditivity, and symmetry is immediate. For $(a,b)\in O$ the adjoint identity is explicitly
\[
 h_1((aX_1,bX_2),Y)=aX_1Y_2+bX_2Y_1
 =h_1(X,(bY_1,aY_2)).
\]
Take $w_0=(1,0)$, so $w_0+w_0^*=1$ and $h_{1,L}=X_1X_2$. For $f=1$ the normalized series is $\exp(X_1X_2)$, whose coefficients $1/n!$ have unbounded denominators. If $\partial_U\partial_VP=0$, its selected value agrees with that in gauge zero for every $f$, by \eqref{eq:fischer-ideal} with $Q=X_1X_2$. Equivalently $P$ is a sum of pure polynomials in $U$ and in $V$, and the exponential is one on each axis. The second-order operator in the source's symmetric-matrix convention is $2\partial_U\partial_V$, with the same kernel over $\Q_p$, including when $p=2$. For $f=1$ all positive-degree selected values vanish. This proves (a) and shows why primitive initial terms and full-series nonintegrality do not by themselves decide harmonic integrality.

For (b), use DP38's invariant splitting in logarithm coordinates, whose coefficient is $[x^py_1]B=-1/p$. On two copies put
\[
 \Theta=B(x_1,\mathbf y^{(1)})^{1/2}B(x_2,\mathbf y^{(2)})^{1/2},\qquad
 P=X_1^pY_{1,1}-pX_1^{p-1}X_2Y_{1,2}.
\]
Give the positive type one row and the negative type $d$ rows, each with two columns. Give the positive row degree $p$, the first negative row degree one, and all other rows degree zero. Every coefficient is integral and every monomial has those degrees. The positive row has no earlier row. Adding an earlier negative row to a later one leaves $P$ unchanged because it uses only the first negative row. Thus the prescribed row invariances hold. For the identity-form contractions
\[
 \Delta_i=\partial_{X_1}\partial_{Y_{i,1}}+
          \partial_{X_2}\partial_{Y_{i,2}},
\]
one has $\Delta_iP=0$ for $i>1$ and
\[
 \Delta_1P=pX_1^{p-1}-pX_1^{p-1}=0.
\]
These are exactly the algebraic definitions \path{lap}, \path{harmK}, \path{RowInv}, \path{RowDeg}, and \path{HaO} in
\path{Formalism/LocalToGlobal/ProjectiveLattices/RowDeg.lean},
lines 21--30 and 37--48; \path{IsThRule.lean}, \path{harmA}, lines 86--87, specializes this definition in the source. They are sums over columns, not individual contractions.

Every monomial of $B-1$ has positive degree on both sides. Therefore the $y_1$-linear coefficient of $B^{1/2}$ is half that of $B$, and the first monomial of $P$ contributes
\[
 p!\,[x^py_1]B^{1/2}=-\frac{(p-1)!}{2},
\]
a unit. The second monomial contributes zero: it differentiates the first copy $p-1>0$ times in $x_1$, with its entire dual variable set zero, where that factor is identically one. Hence $P(\partial)\Theta(0)=-(p-1)!/2$.

For an integral function $f$ on $(G_h\times H)^2$, the first term of $P(D)f$ is divisible by $p$ because $D_{x_1}^p$ is coefficientwise divisible by $p$ at height at least two. Its other invariant derivation preserves that ideal. The second term has its explicit factor $p$ and otherwise only integral invariant derivations. Thus $P(D)f(0)\in pR$. This proves (b), with no assumption that the formal row labels already define a new global source datum.
\end{proof}

The associated ideal in \eqref{eq:fischer-ideal} is generated by column sums $\sum_bX_{\tau,a,b}X_{\bar\tau,a',b}$, not every separate product. In DP22's pure CM type $\{1,2,3\}\subset(\Z/7)^\times$, the conjugate type is $\{6,5,4\}$ and is disjoint. If this formula is applied with one row on the type and zero outside it, each conjugate-row contraction has an empty index set. It does not impose $\partial_x\partial_{y_i}$ on the height-three line and its dual directions. This observation is about the stated harmonic formula; DP25 still prevents treating that row list as the old induced-type global host.

